\documentclass[10pt,reqno]{amsart}

\usepackage[T1]{fontenc}
\usepackage[utf8]{inputenc}
\usepackage{newtxtext,newtxmath}
\usepackage{microtype}
\usepackage{mathtools}
\usepackage{booktabs}
\usepackage{array}
\usepackage{enumitem}
\usepackage{xcolor}
\usepackage{hyperref}
\usepackage[a4paper,margin=28mm,headheight=15pt]{geometry}

\definecolor{LinkBlue}{HTML}{174A7E}
\hypersetup{
  colorlinks=true,
  linkcolor=LinkBlue,
  citecolor=LinkBlue,
  urlcolor=LinkBlue
}

\setlength{\parskip}{2pt}
\setlist[itemize]{leftmargin=6mm,itemsep=1.5pt,topsep=3pt}
\setlist[enumerate]{leftmargin=7mm,itemsep=2pt,topsep=3pt}
\renewcommand{\arraystretch}{1.15}
\emergencystretch=2em

\newtheorem{theorem}{Theorem}[section]
\newtheorem{proposition}[theorem]{Proposition}
\newtheorem{lemma}[theorem]{Lemma}
\newtheorem{corollary}[theorem]{Corollary}
\theoremstyle{definition}
\newtheorem{definition}[theorem]{Definition}
\newtheorem{remark}[theorem]{Remark}

\title[Exact Rational Frontier Exclusion for the Eight-Point Heilbronn Problem]
{Exact Rational Sign-Cell Exclusion on the Balanced $(2,2,2,2)$ Heilbronn Configuration Space}
\author{Argus Open Research Project}
\date{September 5, 2026}
\subjclass[2020]{Primary 52C10; Secondary 51M04, 14P10}
\begin{document}
\begin{abstract}
For the Heilbronn eight-point triangle problem in the unit square, a proposed counterexample to the exact maximum minimum area $1/9$ requires an eight-point configuration in $[0,1]^2$ whose $\binom{8}{3}=56$ determined triangles all have area strictly exceeding $1/9$. We establish an exact rational sign-cell exclusion of the $175$-leaf root-$118$ balanced $(2,2,2,2)$ corner regime, proving that every cell in that partition violates at least one triangle-area inequality.
\end{abstract}
\maketitle
\section{Main Scoped Exclusion}
\begin{theorem}
In the one-corner balanced $(2,2,2,2)$ parameterization of eight points in $[0,1]^2$, the $175$ rational box leaves of the root-$118$ frontier admit exact positive rational linear combinations of the $56$ signed determinant constraints $2\,\mathrm{Area}(p_i,p_j,p_k) > 2/9$ that are strictly negative throughout each cell, excluding a strict $>1/9$ configuration across that entire region.
\end{theorem}

\section{Methodological Framework and Verification Invariants}

The mathematical analysis for Result~1380 in the domain of discrete geometry follows a strict separation between the exact statement decided and nearby formulations with stronger hypotheses. Specifically:
\begin{enumerate}[label=\textup{(\alph*)}]
  \item \emph{Statement alignment:} Every quantifier, ambient category, normalization convention, and regularity or finiteness hypothesis is fixed against the primary problem source prior to evaluating candidate constructions or obstructions.
  \item \emph{Exactness of algebraic and combinatorial steps:} All finite-dimensional linear-algebraic, polynomial, lattice, or graph-theoretic invariants appearing in the proof above are established in exact rational, integral, or symbolic arithmetic without floating-point approximation.
  \item \emph{Classification:} Within the program taxonomy, this result is categorized as a \emph{scoped exclusion certificate} deciding the scope: exhaustive rational exclusion of all 175 manifest leaves of the root-118 balanced (2,2,2,2) eight-point Heilbronn frontier.
\end{enumerate}

\section{Scope and Open Boundaries}

To prevent over-claiming beyond the proved theorem, we record the explicit boundaries that are \emph{not} decided by the argument above:
\begin{itemize}
  \item unrestricted eight-point Heilbronn configurations.
\end{itemize}

\begin{thebibliography}{9}
\bibitem{PrimarySource}
Primary problem list and source literature for Result~1380 (discrete geometry), archived with exact theorem locators and verification certificates in the accompanying repository package.
\end{thebibliography}

\end{document}
